\section{Introducción}

La movilidad urbana representa uno de los principales desafíos en las grandes ciudades debido al crecimiento poblacional y al incremento constante de los medios de transporte. Como alternativa sustentable, los sistemas de bicicletas compartidas permiten a los usuarios desplazarse de manera eficiente, económica y amigable con el medio ambiente.

El presente proyecto consiste en el diseño e implementación de una base de datos para el sistema ECOBICI, cuyo objetivo es administrar la información relacionada con usuarios, membresías, estaciones, bicicletas, viajes, incidentes, métodos de pago y personal operativo. La base de datos fue desarrollada aplicando los principios del diseño conceptual, lógico y físico, garantizando la integridad, consistencia y disponibilidad de la información.

Asimismo, se implementaron mecanismos de seguridad, restricciones de integridad, procedimientos almacenados, funciones, vistas y disparadores que permiten automatizar procesos y asegurar el correcto funcionamiento del sistema. Finalmente, se desarrollaron consultas e informes estadísticos que facilitan la toma de decisiones dentro de la organización.
